% ===== PRÉAMBULE CANONIQUE — à coller VERBATIM en tête de CHAQUE .tex =====
\documentclass[11pt,a4paper]{article}
\usepackage[utf8]{inputenc}\usepackage[T1]{fontenc}\usepackage{lmodern}
\usepackage[french]{babel}
\usepackage{amsmath,amssymb,mathtools}\usepackage{icomma}
\usepackage{siunitx}\sisetup{locale=FR}  % \qty{}{}, \si{}, \ang{}
\usepackage{xcolor}\usepackage{colortbl}
\usepackage{tikz}\usepackage{pgfplots}\pgfplotsset{compat=1.18}
\usepackage[siunitx,europeanresistors]{circuitikz} % circuits (élec.)
\usepackage[most]{tcolorbox}\usepackage{enumitem}\usepackage{array,booktabs}
\usepackage[a4paper,margin=2.2cm]{geometry}
\usetikzlibrary{arrows.meta,calc,angles,quotes,positioning,patterns,%
  backgrounds,shapes.geometric,decorations.pathmorphing,decorations.markings,intersections}
\definecolor{primary}{RGB}{222,49,99}\definecolor{primarydark}{RGB}{150,22,55}
\definecolor{defcol}{RGB}{0,114,178}\definecolor{methcol}{RGB}{52,92,148}
\definecolor{excol}{RGB}{0,135,90}\definecolor{attncol}{RGB}{200,40,55}
\tcbset{boxbase/.style={enhanced,breakable,boxrule=0.9pt,arc=2.5pt,
  left=3mm,right=3mm,top=2.3mm,bottom=2.3mm,fonttitle=\bfseries,coltitle=white}}
% --- boîtes TUTORIEL (noms EXACTS, ne pas renommer) ---
\newtcolorbox{cle}[1][]{boxbase,colback=defcol!6,colframe=defcol,
  title={\textsc{À retenir}\ifx\relax#1\relax\relax\else\textmd{ : \itshape #1}\fi}}
\newtcolorbox{methode}[1][]{boxbase,colback=methcol!7,colframe=methcol,sharp corners,
  title={\textsc{Méthode}\ifx\relax#1\relax\relax\else\textmd{ --- \itshape #1}\fi}}
\newtcolorbox{exemplebox}[1][]{boxbase,colback=excol!5,colframe=excol,
  title={\textsc{Exemple}\ifx\relax#1\relax\relax\else\textmd{ : \itshape #1}\fi}}
\newtcolorbox{piege}[1][]{boxbase,colback=attncol!6,colframe=attncol,
  title={\textsc{Piège}\ifx\relax#1\relax\relax\else\textmd{ : \itshape #1}\fi}}
\newtcolorbox{remarque}[1][]{boxbase,colback=black!4,colframe=black!45,
  title={\textsc{Remarque}\ifx\relax#1\relax\relax\else\textmd{ : \itshape #1}\fi}}
% --- boîtes EXERCICE / CORRIGÉ (noms EXACTS, auto-comptées) ---
\newtcolorbox[auto counter]{exo}[1][]{boxbase,colback=primary!4,colframe=primary,
  title={\textsc{Exercice}~\thetcbcounter\ifx\relax#1\relax\relax\else\textmd{ --- \itshape #1}\fi}}
\newtcolorbox[auto counter]{corrige}[1][]{boxbase,colback=excol!4,colframe=excol,
  title={\textsc{Corrigé}~\thetcbcounter\ifx\relax#1\relax\relax\else\textmd{ --- \itshape #1}\fi}}
\newcommand{\vv}[1]{\vec{#1}}\newcommand{\ds}{\displaystyle}
\newcommand{\R}{\mathbb{R}}\newcommand{\N}{\mathbb{N}}\newcommand{\Z}{\mathbb{Z}}
% (un \usepackage{fancyhdr}+en-tête est possible ; il est ignoré côté web)
% =========================================================================
\begin{document}
\sisetup{per-mode=symbol}

\begin{corrige}[à quelle altitude le skieur atteint-il sa vitesse ?]
\begin{enumerate}[label=\alph*)]
  \item $v=\dfrac{72}{3,6}=\qty{20}{\metre\per\second}$. Conservation : $\tfrac12 v^2=g\,\Delta h$, donc
        $\Delta h=\dfrac{v^2}{2g}=\dfrac{20^2}{2\times10}=\dfrac{400}{20}=\qty{20}{\metre}$.
  \item Le skieur a perdu $\qty{20}{\metre}$ d'altitude : il est à
        $2500-20=\qty{2480}{\metre}$.
\end{enumerate}
\end{corrige}

\begin{corrige}[vitesse après une descente sur plan incliné]
\begin{enumerate}[label=\alph*)]
  \item $h=L\sin\ang{30}=10\times0,5=\qty{5}{\metre}$.
  \item $v=\sqrt{2gh}=\sqrt{2\times10\times5}=\sqrt{100}=\qty{10}{\metre\per\second}$.
\end{enumerate}
\end{corrige}

\begin{corrige}[la balle qui rebondit]
\begin{enumerate}[label=\alph*)]
  \item $\Delta h=\dfrac{v^2}{2g}=\dfrac{10^2}{2\times10}=\dfrac{100}{20}=\qty{5}{\metre}$ au-dessus de
        la plateforme.
  \item Hauteur maximale par rapport au sol : $3+5=\qty{8}{\metre}$.
\end{enumerate}
\end{corrige}

\begin{corrige}[le sac lâché d'une montgolfière]
\begin{enumerate}[label=\alph*)]
  \item Au lâcher, le sac possède la vitesse de la montgolfière : $\qty{10}{\metre\per\second}$
        \emph{vers le haut}. Il continue donc de monter de
        $\Delta h=\dfrac{v^2}{2g}=\dfrac{10^2}{2\times10}=\qty{5}{\metre}$ avant de redescendre.
  \item Hauteur maximale par rapport au sol : $120+5=\qty{125}{\metre}$.
\end{enumerate}
\end{corrige}

\begin{corrige}[le wagon des montagnes russes]
\begin{enumerate}[label=\alph*)]
  \item Conservation entre $A$ et $B$ : $\tfrac12 m v_B^2=mg\,h_A$, d'où
        $v_B=\sqrt{2g\,h_A}=\sqrt{2\times10\times45}=\sqrt{900}=\qty{30}{\metre\per\second}$.
  \item Entre $A$ et $C$ : $\tfrac12 m v_C^2=mg(h_A-h_C)$, donc
        $v_C=\sqrt{2g(h_A-h_C)}=\sqrt{2\times10\times25}=\sqrt{500}\approx\qty{22,4}{\metre\per\second}$.
\end{enumerate}
\end{corrige}
\end{document}
